% Solutions/hints for ch08 exercises (assembled into Appendix C)

\begin{solution}{exr:ch08-table}
(a) Agent~2 takes the straight path $\pi_2[t] = (0,t)$: $R_V$ receives $((0,t), t)$ for $t = 0, \dots, 6$, $R_E$ receives $(((0,t),(0,t{+}1)), t)$ for $t = 0, \dots, 5$, and $(0,6)$ is parked from $t = 6$. Agent~3 then finds nothing in its way and goes straight down, $(0,3), (1,3), (2,3)$: $R_V$ gains $((0,3),0)$, $((1,3),1)$, $((2,3),2)$, $R_E$ gains $(((0,3),(1,3)),0)$ and $(((1,3),(2,3)),1)$, and $(2,3)$ is parked from $t = 2$. (b) The goal $(2,6)$ of agent~1 is never reserved, so $t_g = -1$: any arrival time is acceptable. (c) From $(2,0)$ agent~1 can reach $(2,1)$ and $(2,2)$, but $(2,3)$ is parked from $t = 2$ and agent~1 cannot be there before $t = 3$, so the bottom corridor is cut. The other route climbs through $(1,0)$ to $(0,0)$ at $t = 2$ (free: agent~2 left it at $t = 0$) and follows agent~2 along row~$0$ two steps behind it, which following conflicts allow, but $(0,6)$ is parked from $t = 6$ and agent~1 would arrive there at $t = 8$. Since $(1,6)$ and the cells $(2,4), (2,5), (2,6)$ are reachable only through $(0,6)$ or $(2,3)$, the goal is unreachable. The implementation collapses the time index beyond $t_{\mathrm{static}} = H + 1 = 7$, the last timed entry being agent~2's arrival at $t = 6$; every reachable cell is therefore expanded at most eight times, the open list runs empty, and the function returns \code{None}.
\end{solution}

\begin{solution}{exr:ch08-following}
(a) Reserve every occupied cell for one step longer: for a path with arrival time $T$ put both $(\pi[t], t)$ and $(\pi[t], t{+}1)$ into $R_V$ for $0 \le t \le T$ (equivalently, test $(v', t)$ as well as $(v', t{+}1)$ before entering $v'$; the parked goal already covers all later times). A cell that a higher-priority agent leaves at $t$ is then blocked at $t+1$, which is exactly the one-cell gap. The swap test becomes redundant under this convention: to exchange cells with an agent that moves $u \to v$ during step $t$, the lower-priority agent would have to enter $u$ at $t+1$, and the trailing entry $(u, t{+}1)$ already forbids that. Keep the test all the same -- it costs one lookup, and it keeps the two conventions independent, so that switching the gap off does not silently allow swaps. (b) Under disappear-at-target an agent vanishes when it arrives, so no goal is parked on and the goal test of line~\ref{alg:ch08-ca:goal} loses its condition $t > t_g$. Agent~3 then reserves only $((0,3),0)$, $((1,3),1)$, $((2,3),2)$ and the two moves between them, and it is gone from $t = 2$ on. Agent~1 crosses $(2,3)$ at $t = 3$ and agent~2 crosses $(0,3)$ at $t = 3$, one step after agent~3 has left, which is a following move and therefore allowed. No agent ever blocks another, so all six orders succeed with the independent shortest paths, costs $6, 6, 2$ and $\sumcost = 14$: exactly the lower bound quoted under \cref{tab:ch08-orders}. The whole difficulty of \cref{ex:ch08-three} lives in the parked reservations, and this is why the convention must be stated before any claim about prioritized planning is made.
\end{solution}

\begin{solution}{exr:ch08-pocket}
(a) With $\pi_1[t] = (1,t)$ for $t \le 4$ and $(1,4)$ parked afterwards, agent~2 starts at $(1,4)$. At $t = 1$ it can be at $(1,4)$ or $(1,3)$. At $t = 2$ the cells $(1,4)$ and $(1,3)$ remain; $(1,2)$ is agent~1's cell at that time. At $t = 3$ agent~1 is at $(1,3)$, so agent~2 must be at $(1,4)$, which it reaches by waiting or by moving back from $(1,3)$ (not a swap, because agent~1 comes from $(1,2)$). At $t = 4$ agent~1 arrives at $(1,4)$: waiting is a vertex conflict and moving to $(1,3)$ exchanges cells with agent~1, a swap. The set of legal states is empty and the search fails; the pocket $(0,2)$ was never reachable, since it requires being at $(1,2)$ at $t = 1$, two cells away from the start. The order $(2,1)$ is the mirror image. (b) The second pocket $(0,3)$ is adjacent to $(0,2)$, so the two pockets form a two-cell bypass above $(1,2)$ and $(1,3)$. The order $(1,2)$ now succeeds: agent~1 goes straight (cost $4$) and agent~2 takes $(1,4), (1,3), (0,3), (0,2), (1,2), (1,1), (1,0)$, climbing into the bypass at $t = 2$ while agent~1 is at $(1,2)$, moving along it at $t = 3$ while agent~1 is at $(1,3)$, and coming down at $t = 4$ as agent~1 parks at $(1,4)$; cost $6$, sum $10$. The order $(2,1)$ still fails: agent~2's straight path reaches $(1,2)$ at $t = 2$ and $(1,1)$ at $t = 3$, and agent~1, two cells from the nearest pocket, cannot get out of the corridor in time. (c) Yes. The lower bound is $4 + 4 = 8$, and two agents cannot pass each other inside a corridor, so one of them must leave it, which costs at least two extra steps (one up, one down); the plan of (b) spends exactly these two steps and nothing else, so $10$ is optimal, as \code{optimal\_soc} confirms.
\end{solution}

\begin{solution}{exr:ch08-wellformed}
(a) For each agent $i$ remove the $2k - 2$ other endpoints from $G$ and run one breadth-first search from $s_i$; the instance is well-formed if and only if each search reaches $g_i$. Each search costs $\bigO{\abs{V} + \abs{E}}$, hence $\bigO{k(\abs{V} + \abs{E})}$ in total. (b) The step ``waiting is unblocked'': it needs that no entry $(s_i, t)$ exists in $R_V$, which holds only because every earlier agent had $s_i$ in its static obstacle set. For the instance take a $3 \times 5$ open grid with agent~1 from $(0,0)$ to $(0,4)$ and agent~2 from $(0,2)$ to $(1,2)$. It is well-formed: agent~1 can avoid both endpoints of agent~2 through the bottom row, and agent~2 reaches its goal in one step. Without the rule, agent~1's shortest path runs along the top row and reserves $(s_2, 2) = ((0,2), 2)$; agent~2 must have left its start by then, which here it can do by moving to its goal at $t = 1$, but nothing in the algorithm guarantees such an escape in general. With the rule agent~1 treats $(0,2)$ as an obstacle and takes a path of length $6$ around it, through $(1,1), (1,2), (1,3)$; agent~2 then has to step back to $(0,2)$ while agent~1 passes its goal and arrives at $t = 4$, so the guaranteed plan costs $10$ where the unguarded one cost $5$. (c) On a $3 \times 3$ open grid let agent~1 go from $(1,0)$ to $(1,2)$ and agent~2 from $(0,1)$ to $(1,1)$. Both endpoint-free paths exist (agent~1 around the bottom row, agent~2 straight down), so the instance is well-formed, but agent~1's shortest path passes through $g_2 = (1,1)$. In the order $(2,1)$ agent~2 parks on $(1,1)$ at $t = 1$ and \cref{alg:ch08-ca} returns a detour of cost $4$ for agent~1 through row~$0$ or row~$2$; the code takes $(1,0), (0,0), (0,1), (0,2), (1,2)$, passing through $s_2$ after agent~2 has left it. In the order $(1,2)$ agent~1 takes the straight path and visits $(1,1)$ at $t = 1$, so $t_g = 1$ for agent~2, which waits one step and arrives at $t = 2$; both orders succeed, as the theorem promises.
\end{solution}

\begin{solution}{exr:ch08-whca}
(a) Let the agents be adjacent, agent~1 at $(1,3)$ and agent~2 at $(1,4)$, and let agent~1 be planned first in this round. The table starts empty, so agent~1 does not see agent~2 at all: it plans straight ahead, $(1,3), (1,4), (1,5), \dots$, and reserves those cells. Agent~2 now finds every cell in front of it taken: stepping to $(1,3)$ at $t = 1$ would swap with agent~1's move $(1,3) \to (1,4)$, and waiting at $(1,4)$ collides with agent~1 there at $t = 1$. Its cheapest unblocked path is to retreat, $(1,4), (1,5), (1,6), \dots$, and to turn around later -- beyond the window, where agent~1 is invisible. Each agent executes one step: agent~1 advances to $(1,4)$, agent~2 falls back to $(1,5)$. Line~\ref{alg:ch08-whca:rotate} then makes agent~2 the first planner, the roles are exchanged, and agent~1 is pushed back to $(1,3)$. The pair oscillates between the two configurations for ever.

(b) The window size is irrelevant because the agent that plans first never sees the other one: at the start of every round the table is empty, and its own path is a straight line through the other agent. Only the second planner sees a conflict, and its only escape is backwards. Enlarging $w$ lets the second planner look further ahead -- it plans a longer retreat and a turn -- but only the first $\delta = 1$ step of that plan is executed, and the retreat is the first step. The single pocket $(0,1)$ sits at the far end of the corridor, behind agent~1's start, so no window that reaches it from the meeting point is ever used: the agents meet near $(1,4)$ and never get back. \code{whca\_star} therefore reaches \code{max\_rounds} and returns \code{None} for every $w = 2, \dots, 12$.

(c) Rotation is what causes the symmetry: it hands the right of way to the two agents alternately, so neither ever completes a passing manoeuvre. Freezing the order for several rounds is what helps. With the full horizon and the fixed order $(2,1)$ the instance is solved: agent~2 goes straight and arrives at $t = 8$, and agent~1 first retreats into the pocket $(0,1)$, lets agent~2 pass and then walks the whole corridor, arriving at $t = 15$ (\code{prioritized\_planning} returns costs $15$ and $8$). Any of the following removes the deadlock: keep the priority order fixed until every agent has reached its goal; grow the window until each agent can see a pocket, and re-plan with the same order; or detect the repeated configuration and hand the two agents to a coupled solver (\cref{ch:ch09}) for the corridor.
\end{solution}
